\documentclass[11pt]{article}
\title{On the Closure of Compositional Hierarchies in Digital Symmetries}
\author{A rigorous researcher-engineer bridging abstract mathematics and concrete digital systems.}
\usepackage[margin=1in]{geometry}
\usepackage{graphicx}
\usepackage{amsmath,amssymb,mathtools}
\usepackage{hyperref}
\graphicspath{{media/}{media/diagrams/}{images/}{artwork/}}

\begin{document}

\section*{Glossary \& Symbol Index}

This back-matter section serves two purposes: it provides plain-language definitions for recurring terms, and it gives a bidirectional map from symbols to the chapters and sections where they are introduced, used, or proved. The intent is to make the book navigable for readers who want to move from notation to meaning, and from a symbol back to the relevant proof, construction, or build artifact.

\begin{description}
  \item[Glossary.] Concise definitions of the book's recurring technical terms are given in reader-facing language. Each entry should point, where appropriate, to the chapter, section, proposition, theorem, or build artifact in which the term is first defined or most substantially developed.
  \item[Symbol index.] A symbol-to-location map records where each notation item appears in the book. The map is bidirectional in the sense that readers can locate a symbol from its meaning, and can also recover the symbol from the chapter or section in which it is used.
  \item[Cross-references.] Definitions and symbol entries should cross-reference proofs, constructions, and worked examples rather than restating them. This section is therefore a navigational layer, not a substitute for the main exposition.
\end{description}

\medskip

\noindent\textbf{Glossary entries.} The following terms are representative of the vocabulary used throughout the book; additional entries may be added as the manuscript stabilizes.

\begin{description}
  \item[Closure.] A property of a class of objects or operations indicating that the class is stable under a specified operation. In this book, closure is used in the sense of algebraic, categorical, and computational stability conditions; see the chapters on clones, traced composition, and symmetry-aware analysis for the main uses.
  \item[Composition.] The operation of combining processes, maps, circuits, or morphisms so that the output of one becomes the input of another. Composition is the organizing principle of the book's hierarchy of constructions, from basic Boolean maps to modular and traced systems.
  \item[Symmetry.] A transformation that preserves the relevant structure of an object, function, circuit, or system. Symmetry may appear as an automorphism, an invariance, an equivariance condition, or a group action, depending on the chapter.
  \item[Clone.] A set of Boolean functions closed under composition and containing projections. Clones are central to the discussion of Post's lattice and functional completeness.
  \item[Operad.] An algebraic structure encoding multi-input composition with typed arities. Operads provide a language for compositional hierarchies and modular assembly.
  \item[PROP.] A symmetric monoidal category-like formalism for operations with multiple inputs and outputs. PROPs are used to model circuit-like and process-like composition.
  \item[Trace.] A feedback operator that closes a wire or channel to form a loop. Traces are used to formalize sequential composition with feedback.
  \item[Equivariance.] Compatibility of a map or model with a symmetry action. Equivariant architectures preserve symmetry information by construction.
  \item[Reproducibility.] The property that a claim can be rebuilt, checked, or rerun from the stated artifacts, scripts, and inputs. In this book, reproducibility is treated as part of the mathematical workflow.
\end{description}

\medskip

\noindent\textbf{Symbol index.} The table below is a placeholder structure for the bidirectional symbol map. It should be populated with the final manuscript's notation and the exact chapter, section, and page locations once pagination is fixed.

\begin{center}
\begin{tabular}{|l|l|l|}
\hline
\textbf{Symbol} & \textbf{Meaning} & \textbf{Location(s)} \\
\hline
$\mathbb{B}$ & Boolean domain & Chapter~3; Chapter~7; Chapter~14 \\
\hline
$\mathbb{B}^n$ & $n$-fold Boolean cube & Chapter~1; Chapter~14; Chapter~15 \\
\hline
$\circ$ & Composition of maps or morphisms & Chapters~1--16 \\
\hline
$\otimes$ & Tensor or parallel composition & Chapters~2, 5, 12 \\
\hline
$\mathrm{Cl}(\cdot)$ & Closure operator or closure of a set & Chapter~3; Chapter~15 \\
\hline
$\mathrm{Sym}(\cdot)$ & Symmetry group or symmetry set & Chapter~8; Chapter~13 \\
\hline
$\mathrm{Tr}$ & Trace / feedback operator & Chapter~5; Appendix material on traced monoidal structure \\
\hline
$\mathrm{NPN}$ & Negation-permutation-negation equivalence & Chapter~8; Chapter~15 \\
\hline
\end{tabular}
\end{center}

\medskip

\noindent\textbf{Usage note.} For each symbol, the index should list the earliest defining location, the principal development location, and any later sections where the notation is reused in a materially different context. When page numbers are available, they should be included alongside chapter and section references to support fast lookup in the printed edition.

\medskip

\noindent\textbf{Editorial convention.} If a term has both an informal and a formal meaning in different chapters, the glossary entry should state the informal meaning first and then point to the formal definition. If a symbol is overloaded, the index should disambiguate it by context rather than by inventing new notation unless the manuscript explicitly introduces one.

\medskip

\noindent\textbf{Revision note.} To align this section more tightly with the rest of the manuscript, the glossary should remain concise and reader-facing, while the symbol index should be expanded only after the notation in Chapters~1--16 is finalized. In particular, the final pass should add section-level and page-level references for the most frequently reused symbols, and should distinguish between definitions, principal uses, and proof locations where those differ.


\end{document}
